\section{OpenAI、GPT-6 Sol / Luna を投入}
OpenAI公式はGPT-6系列にSolとLunaを追加した。Solは作業・コーディング・エージェント、Lunaは安価・高速帯を担うと説明し、API価格をGPT-5.6のプロモ価格比で50\%低減したとする。公式ブログではSolを入力\$2・出力\$10、Lunaを入力\$0.10・出力\$0.50／100万トークンと掲載した。
\dailyxsource{OpenAI (@OpenAI)}{https://x.com/OpenAI/status/2102460975790137662}

\section{Anthropic、Claude Opus 5.5 を公開}
Anthropic / Claude公式はClaude 5.5ファミリー最初のモデルとしてOpus 5.5を公開した。Opus 5より典型ワークロードで約40\%安く、Fast modeは最大約2.5倍と説明。System cardと外部評価も同日前後に提示されたが、ベンチマーク値は各資料の条件に依存する。
\dailyxsource{Claude (@claudeai)}{https://x.com/claudeai/status/2102435511222890900}

\section{同日リリースを「価格戦争」と読む反応}
観測窓終盤には、OpenAIとAnthropicが同じ夜に低コスト化を前面に出したことから「価格戦争」と整理する投稿が現れた。両社の値下げ・新モデル投入は公式情報だが、協調や戦略的意図まで一次資料で確認されたものではない。
\dailyxsource{Ayush (@chillashw)}{https://x.com/chillashw/status/2102518272813879497}
\dailyxnote{「価格戦争」は観測者の解釈として扱う。}

\section{Muse、PayPal・Expedia連携とSpark 1.3を訴求}
Alexandr WangはPayPalでのMuseチェックアウト連携とExpediaの旅行連携を歓迎し、connector platformの開放にも言及した。Scale側はSpark 1.3とPRBenchでの金融・法律領域首位を主張する一方、ダウンロード数や株価影響などは二次報道・投稿者主張が混在している。
\dailyxsource{Alexandr Wang (@alexandr\_wang)}{https://x.com/alexandr_wang/status/2102418455622508722}

\section{AmazonがMuseの買い物を遮断したとの報道}
窓内では、AmazonがMuseを無断エージェントとして買い物から排除したとの報道が再流通した。ブロック実施自体は9月20日前後とされ観測窓外。Amazon側の規約・識別・資格情報への懸念と、Meta側の安全設計とされる説明はいずれも二次要約が中心で、公式声明本文はGrok収集では未取得だった。
\dailyxsource{TNW (@thenextweb)}{https://x.com/thenextweb/status/2102509680236822723}
\dailyxnote{イベント自体は観測窓外で、窓内の再流通を収録。}

\section{Grok Botデスクトップ、53件の性能修正}
開発者laurenはGrok Botデスクトップへ直近数日で53件の性能修正を入れたと報告した。再接続60秒から0.7秒、スリープ復帰23秒から1秒、Computer view切替868msから27msなどの改善値を提示している。数値は開発者自己報告で独立計測ではない。
\dailyxsource{lauren (@poteto)}{https://x.com/poteto/status/2102504221648339170}

\section{Tesla車内Grok、Connectorsで外部作業へ}
Tesla公式は、車内GrokがConnectorsを通じて受信箱、カレンダー、既存ファイル・チャット・タスクをハンズフリーで扱えるようになったと発表した。対象車両・地域・対応コネクタの詳細は投稿本文だけでは不明で、プライバシー面の一次資料も今後の確認事項として残る。
\dailyxsource{Tesla (@Tesla)}{https://x.com/Tesla/status/2102430656349544590}

\section{TypeSafe AIのJev、需要急増で登録を一時停止}
TypeSafe AIは、テキストを生成せず判定・選択・スコアを返すdecision model「Jev」への需要急増を受け、新規サインアップを一時停止した。MotherDuckやPydanticとの統合も話題となり、LLMの前段で判定やルーティングを担わせる用途が議論された。
\dailyxsource{TypeSafe AI (@typesafeai)}{https://x.com/typesafeai/status/2102281508950307159}

\section{OpenAI、第三者評価の優先領域と原則を公開}
OpenAIは、訓練・評価・デプロイへの深いアクセスを伴う独立評価を支援すると表明し、安全ケース、セーフガード、Preparedness系能力評価などの優先領域と実施原則を整理した。制度的な強制監査そのものではなく、実際の評価機関と適用モデルが今後の焦点となる。
\dailyxsource{OpenAI (@OpenAI)}{https://x.com/OpenAI/status/2102447425243828347}

\section{エージェント群とボット検証需要をめぐる主張}
Nikita Bierは、エージェントスウォームが中小サイトや行政フォームを圧迫し得るとして、ボット検出と人間検証の需要拡大を主張した。X在籍時の経験を基にした市場観察で、定量データや製品発表を伴うものではない。
\dailyxsource{Nikita Bier (@nikitabier)}{https://x.com/nikitabier/status/2102432368158245252}
\dailyxnote{Grok intakeではUnverified。経験談・市場観察として扱う。}

\section{ForecastBench、LLMと予測専門家が有意差なしと報告}
Forecasting Research InstituteはForecastBenchで、LLMモデル群がスーパーフォアキャスターと「統計的有意差なし」で並んだと報告し、10月の再戦を告知した。使用モデルや質問セットの詳細は投稿だけでは不完全で、独立再現も未確認である。
\dailyxsource{Forecasting Research Institute (@Research\_FRI)}{https://x.com/Research_FRI/status/2102411966103093386}

\section{NVIDIA、AI factoryの耐久性を強調}
NVIDIA公式はG20 Innovation MinisterialでのJensen Huangの発言として、GW級AI factoryは500--600億ドル規模になり得るため、モデルやアルゴリズムの変化で陳腐化しにくいfungible / durableな構成が必要だと紹介した。金額はNVIDIA投稿上の説明である。
\dailyxsource{NVIDIA (@nvidia)}{https://x.com/nvidia/status/2102375573192413569}

\section{LLM生成論文の扱いをめぐる禁止論}
Yarin Galは、透かし検出に失敗した論文をdesk rejectし、arXivではLLM執筆論文を禁止すべきだと主張した。投稿爆発や低品質生成物への対応をめぐる規範的意見であり、ICMLやarXivの公式ポリシー変更を示すものではない。
\dailyxsource{Yarin Gal (@yaringal)}{https://x.com/yaringal/status/2102493445483020322}
\dailyxnote{Grok intakeではUnverifiedの意見投稿。}

\section{Muse Spark 1.3 とPRBenchの首位主張}
Scale AI LabsのAakash SabharwalはMuse Spark 1.3について、自社が評価・学習データ・テストに関与したと明記しつつ、PRBenchの金融・法律領域で首位と主張した。評価関与者自身による報告であるため、独立性は限定的とGrok intakeは整理している。
\dailyxsource{Aakash Sabharwal (@aakashsabharwal)}{https://x.com/aakashsabharwal/status/2102485169034944735}
